\documentclass[11pt,twocolumn]{article}

\usepackage[utf8]{inputenc}
\usepackage[T1]{fontenc}
\usepackage{amsmath,amssymb,amsfonts}
\usepackage{graphicx}
\usepackage{booktabs}
\usepackage{hyperref}
\usepackage[margin=0.85in]{geometry}

\title{\textbf{The Darkness Functional: A Diagnostic Framework for Coherence Loss, Interface Mismatch, and Closed-Loop Restoration Under Constraint}}
\author{Gage Fry \\ \small Independent Researcher, Rockford, Illinois, USA \\ \small \texttt{https://github.com}}
\date{November 2026}

\begin{document}

\maketitle

\begin{abstract}
Complex dynamical, informational, and physical systems rarely fail instantaneously; instead, they undergo gradual, silent coherence collapse driven by stochastic noise, boundary degradation, and inter-subsystem mismatch. Standard binary health checks and post-hoc output metrics frequently fail to detect accumulating fragility until macroscopic failure has already occurred. Building on the \textit{Coherence Under Constraint (CUC)} formalism, we introduce the \textbf{Darkness Functional} $D[\Sigma, t]$ and \textbf{Darkness-Aware Persistence} $S(t)$—a lightweight, domain-agnostic diagnostic and closed-loop stabilization architecture. Rather than treating ``darkness'' as an opposing substance or force, we formalize it as a measurable coherence-loss load that exponentially attenuates usable system persistence. We validate this framework across three testbeds: (1) a hierarchical 60-oscillator Kuramoto network subjected to transient perturbation and active phase restoration, (2) a multi-tier AI agent memory benchmark evaluating retrieval fragmentation, cross-store contradiction, and the ``high-throughput illusion,'' and (3) \textbf{CUC-Array v1}, a low-cost 40~kHz ultrasonic phased-array levitator implementing real-time acoustic phase restoration and Navier--Stokes cascade-gate telemetry on embedded hardware.
\end{abstract}

\section{Introduction}
Across coupled physical oscillators, standing-wave acoustic traps, and long-horizon retrieval-augmented artificial intelligence (AI) agents, structural persistence depends on maintaining phase alignment across constrained state manifolds. In multi-scale systems, local subsystems can exhibit high internal order while global alignment deteriorates—a condition in which local coherence masks impending systemic collapse.

To quantify hidden fragility prior to macroscopic failure, we formalize a diagnostic functional that aggregates four distinct modes of structural degradation into a single scalar load, paired with a closed-loop phase-restoration control law.

\section{Mathematical Formulation}

\subsection{The Darkness Functional}
For a constrained system $\Sigma$ evolving over state space $\Omega$ at time $t$, we define the \textbf{Darkness Functional} $D[\Sigma, t] \in \mathbb{R}_{\ge 0}$ as a weighted linear combination of four normalized observables:
\begin{equation}
D[\Sigma, t] = \alpha\bigl(1 - r(t)\bigr) + \beta N(t) + \gamma F_C(t) + \delta M_I(t)
\label{eq:darkness}
\end{equation}
where $\alpha, \beta, \gamma, \delta > 0$ are domain-calibrated sensitivity weights, and:
\begin{itemize}
    \item $r(t) \in [0, 1]$ is the global phase or semantic synchronization order parameter, so $1 - r(t)$ captures global alignment loss.
    \item $N(t) \ge 0$ measures stochastic perturbation load, velocity variance, or off-topic retrieval noise.
    \item $F_C(t) \ge 0$ quantifies boundary degradation or constraint failure outside an admissible manifold $\mathcal{S} \subseteq \Omega$.
    \item $M_I(t) \ge 0$ measures subsystem interface mismatch, capturing both inter-group dispersion and divergence between local cluster coherence and global coherence.
\end{itemize}

\subsection{Darkness-Aware Persistence and Hysteresis}
Usable system persistence $S(t)$ scales proportionally with global alignment $r(t)$, admissible manifold capacity $C(\Omega, t)$, and sustaining energy or computational throughput $\Phi(t)$, while decaying exponentially under accumulated darkness load:
\begin{equation}
S(t) = r(t) \, C(\Omega, t) \, \Phi(t) \, \exp\bigl(-D[\Sigma, t]\bigr)
\label{eq:persistence}
\end{equation}
Alongside instantaneous load $D[\Sigma, t]$, we define the \textbf{cumulative historical darkness load}:
\begin{equation}
\mathcal{H}_D(t) = \int_0^t D[\Sigma, \tau] \, d\tau
\label{eq:hysteresis}
\end{equation}
Equations~\eqref{eq:darkness}--\eqref{eq:hysteresis} yield five testable predictions: (P1) persistence $S(t)$ decays exponentially as $D[\Sigma, t]$ rises; (P2) interface mismatch $M_I(t)$ can dominate failure even when local subsystems remain synchronized; (P3) closed-loop restoration can recover active persistence $S(t)$ without erasing historical load $\mathcal{H}_D(t)$; (P4) local coherence can mask global instability; and (P5) increasing throughput $\Phi(t)$ without phase alignment accelerates incoherence (the \textit{high-throughput illusion}).

\section{Computational Validation}

\subsection{Hierarchical Kuramoto Network Testbed}
We simulate $N = 60$ coupled phase oscillators $\theta_i \in [-\pi, \pi)$ partitioned into $G = 3$ equal subgroups with intrinsic frequencies $\omega_i \sim \mathcal{N}(0, 0.25^2)$, intra-group coupling $K_{\text{local}} = 1.2$, and inter-group coupling $K_{\text{global}} = 0.4$:
\begin{equation}
\dot{\theta}_i = \omega_i + \sum_{j=1}^{N} \frac{K_{ij}}{N} \sin(\theta_j - \theta_i) + u_{\text{restore}}(t) + \eta_i(t)
\end{equation}
The global order parameter $z = \frac{1}{N}\sum_{j=1}^N e^{i\theta_j} = r e^{i\psi}$ and subgroup order parameters $r_g$ define the interface mismatch:
\begin{equation}
M_I(t) = \frac{1}{G}\sum_{g=1}^{G} \bigl|r_g - \bar{r}_g\bigr| + \max\bigl(0, \bar{r}_g - r\bigr)
\end{equation}
Boundary degradation is evaluated outside the admissible phase window $|\theta_i| \le 1.2$~rad via $F_C(t) = \frac{1}{N}\sum_{i=1}^N \max(0, |\theta_i| - 1.2)^2$. During the perturbation window $t \in (10, 18)$, noise intensity increases from $\sigma = 0.08$ to $\sigma = 0.22$, driving a spike in $D[\Sigma, t]$ and collapsing $S(t)$. During the restoration window $t \in (22, 34)$, feedback control $u_{\text{restore}}(t) = k_{\text{restore}} \sin(\psi - \theta_i)$ suppresses $D[\Sigma, t]$ and restores $S(t)$ above critical gain $k_{\text{crit}}$, while $\mathcal{H}_D(t)$ retains a permanent record of the perturbation interval.

\subsection{AI Memory and Retrieval Case Study}
To evaluate Equations~\eqref{eq:darkness}--\eqref{eq:persistence} in informational systems, we model a 40-turn retrieval-augmented agent querying a multi-tier memory store (\texttt{working}, \texttt{episodic}, \texttt{archival}). Here, $r(t)$ measures query-document cosine alignment discounted by stale records; $N(t)$ captures distractor fraction and score dispersion; $F_C(t)$ tracks missing source provenance and chronological inversions; and $M_I(t)$ measures pairwise factual contradiction across retrieved records.

\begin{table}[h]
\centering
\caption{Cross-Domain Mapping of Darkness Functional Observables}
\label{tab:mapping}
\small
\begin{tabular}{@{}lll@{}}
\toprule
\textbf{Term} & \textbf{AI Memory Store} & \textbf{CUC-Array v1 Hardware} \\
\midrule
$1 - r(t)$ & Semantic alignment loss & Standing-wave lock loss \\
$N(t)$ & Distractor \& score variance & Turbulent acoustic jitter \\
$F_C(t)$ & Missing provenance / order & Emitter phase-drift penalty \\
$M_I(t)$ & Cross-tier factual conflict & Multi-channel phase skew \\
$u_{\text{restore}}$ & Memory consolidation pass & Differential PWM phase shift \\
\bottomrule
\end{tabular}
\end{table}

When contradictory updates and unprovenanced entries are injected (turns 10--19), $M_I(t)$ and $F_C(t)$ spike despite moderate cosine similarity. Doubling retrieval volume ($\Phi = 1.8$, turns 20--27) increases raw token throughput while driving $S(t) \to 0$ via $\exp(-D)$. Executing a provenance-verification and contradiction-pruning consolidation pass at turn 28 restores $S(t)$ within two cycles.

\section{Hardware Implementation: CUC-Array v1}
To validate real-time closed-loop control on physical wave dynamics, we designed \textbf{CUC-Array v1}, a benchtop 40~kHz ultrasonic acoustic levitator driven by an ESP32-S3 microcontroller and L293D dual H-bridge drivers. Two opposing 5-emitter arrays (\texttt{TCT40-16T}, $\lambda \approx 8.58$~mm in air) separated by a resonant cavity height of $34.3\text{ mm} = 4\lambda$ establish standing-wave Gor'kov potential wells spaced at $\lambda/2 \approx 4.29$~mm that suspend 2--3~mm expanded polystyrene spheres.

An electret microphone sensor (\texttt{MAX4466}) samples the acoustic cavity at $\sim 66$~kHz, computing standing-wave coherence $r(t)$, turbulent jitter $N(t)$, and a dimensionless Navier--Stokes cascade ratio $\Gamma_{\text{cascade}}(t)$ at a 100~Hz control rate. When external airflow perturbs the cavity such that $D[\Sigma, t] > D_{\text{thresh}}$ or $\Gamma_{\text{cascade}}(t)$ trips the turbulent streaming gate, the firmware applies proportional differential phase shifts across the bottom emitter bank to recenter the trapping potential prior to particle ejection.

\section{Conclusion and Future Work}
The Darkness Functional provides a computationally inexpensive, mathematically unified diagnostic layer capable of anticipating failure across coupled dynamical networks, agentic AI memory stores, and embedded acoustic phased arrays. Future work includes multi-axis 2D holographic phase steering on CUC-Array v2 and empirical benchmarking across long-horizon enterprise retrieval datasets.

\begin{thebibliography}{9}
\bibitem{kuramoto1984}
Y. Kuramoto, \textit{Chemical Oscillations, Waves, and Turbulence}, Springer-Verlag, Berlin, 1984.

\bibitem{gorkov1962}
L. P. Gor'kov, ``On the forces acting on a small particle in an acoustical field in an ideal fluid,'' \textit{Soviet Physics Doklady}, vol. 6, pp. 773--775, 1962.

\bibitem{marzo2015}
